
    \documentclass[12pt]{article}
    \usepackage{amsmath}
    \usepackage{amsthm}
    \usepackage{amsfonts}
    \usepackage{amssymb}
    \usepackage{enumerate}
    \usepackage{xfrac}
    \usepackage{graphicx}
    \usepackage{mdframed}
    \usepackage{multicol}
\usepackage[T1]{fontenc}
\usepackage{lmodern}
    \usepackage{verbatim}
    \usepackage{tikz}
    \usepackage[margin = .8in]{geometry}
    \geometry{letterpaper}
    \linespread{1.2}
    \newcommand{\RR}{\mathbb{R}} 
    \newcommand{\cringe}[1]{ \( \displaystyle #1 \)}
    \newcommand{\NN}{\mathbb{N}}
    \newcommand{\ZZ}{\mathbb{Z}} 
    \newcommand{\QQ}{\mathbb{Q}}
    \newcommand{\mmod}[1]{\quad (\text{mod }#1)}
    \newcommand\setItemnumber[1]{\setcounter{enumi}{\numexpr#1-1\relax}}
    \newcommand\sgr[1]{\langle #1 \rangle}
    \newcommand\set[1]{\left\lbrace #1 \right\rbrace} 
    \newcommand\abs[1]{\left| #1 \right|} 
    \newcommand\parens[1]{\left( #1 \right)} 
    \newcommand\brac[1]{\left[ #1 \right]} 
    \newcommand\dist[1]{\text{dist}(#1 )}
    \newcommand\ideal[1]{\left\langle #1 \right\rangle}
    
   
    \begin{document}
    


    \noindent Shan Gao \hfill  2023-01-06
    
    \begin{center}
      {\large Math 255 PS1 }
    \end{center}

    \begin{itemize}
        \item[1.] If $G$ is an open set in $\RR$ and $F$ is a closed set in $\RR$, $G \setminus F$ is an open set in $\RR$ and $F\setminus G$ is a closed set in $\RR$. \\
        $G\setminus F = G\cap F^c$ is open since $F$ closed $\implies$ $F^c$ open and open sets are open under finite intersections. $F \setminus G = F \cap G^c $ is closed since $G$ open $\implies$ $G^c$ closed and closed sets are closed under arbitrary intersections. 
        \item[2.]  $S^\circ$ is open since, by definition if $x \in S^\circ$ then $\exists \varepsilon > 0 $ such that $B_\varepsilon(x) \subseteq S$, any $y \in B_\varepsilon(x) \subseteq $ with $y \neq x$ is an interior point, since $B_\varepsilon(x)$ is open, $\exists \varepsilon^\prime > 0 $ such that $B_{\varepsilon^\prime} (y) \subseteq B_\varepsilon(x) \subseteq S$. It is clear that the set of interior point is the union of such open balls. $\displaystyle \bigcup_{x \in S^\circ} B_{\varepsilon_x} (x) = S^\circ$\\
        Let $S$ be an open subset of $\RR$, that means $\forall x \in S$, $\exists \varepsilon > 0$ such that $B_\varepsilon(x) \subseteq S$. The same holds for $x \in S^\circ$. It follows that $x \in S \iff x \in S^\circ$. If $S = S^\circ$, then by the first part of the proof, since $S^\circ$ is open, $S$ is open. 
        \item[3.]
       \begin{proof}
        Consider the subcover, that is the left sided fattening of the set $(0,2)$,  $ O_\alpha =\bigcup_{\alpha \in \NN} \parens{\frac{1}{\alpha}, 2}$, which is the infinite union of open sets of the form $G_\alpha = \parens{\frac{1}{\alpha}, 2}$. This is an open cover of $(0,1]$. Let $x \in (0,1]$ so that $0 < x \leq 1$. It follows that $\exists \varepsilon > 0$ such that $\dist{x,0} \geq \varepsilon$. By the archimedian property of reals, $\exists \alpha_0 \in \NN$ such that  $\frac{1}{\alpha_0} < \frac{\varepsilon}{2}$. So $x$ will be in $G_{\alpha_0}$ since $0 < \frac{1}{\alpha_0} < x \leq 1 < 2$. It follows that $x \in O_\alpha$, so $(0,1] \subseteq O_\alpha$. $O_\alpha$ has no finite subcover in which $(0,1]$ will be entirely contained. Suppose that there exist a finite subcover, that means $(0,1] \subseteq \bigcup_{\alpha \in \set{1,2,\cdots,N}} \parens{\frac{1}{\alpha}, 2}$ for some $N \in \NN$, then since $\frac{1}{N} < \frac{1}{N^\prime}$ if $N > N^\prime$, the largest set in the union is $G_N$, but since $\frac{1}{N} >  0$, consider $\frac{1}{2N}$, which is in $(0,1]$ since $0< \frac{1}{2N} \leq 1$, but it is not in $\bigcup_{\alpha \in \set{1,2,\cdots,N}} \parens{\frac{1}{\alpha}, 2}$ since it is not in any of $G_\alpha$ that compose the union, which leads to a contradiction since it's a subcover of $(0,1]$. 
        \end{proof}
        \item[4.]
        \begin{proof}
        	We prove the corollary stated in class first. Let $K$ be a compact subset of the reals, and $F$ a non-empty closed subset of $K$, then $F \cap K$ is compact. Let $\set{U_\alpha}$ be on open cover of $F$. Since $F$ is closed, It follows that $F^c \cup \set{U_\alpha}$ is an open cover of $K$. But then since $K$ is compact, there exists a finite subcover, call this $U_0 \subseteq F^c \cup \set{U_\alpha}$. But since $F \cap K \subseteq K$, then 
         \begin{align*}
             F \cap K \subseteq K \subseteq U_0
         \end{align*}
         So $U_0$ is a finite cover for $K \cap K$. Since the open cover we chose for $F\cap K$ is arbitrary, we can build a finite subcover for $F \cap K$ for any open cover, thus $F \cap K$ is compact. It follows that if $F$ is a non empty closed subset of $K$, then $F \cap K = F$, by the proof above, there exists a finite subcover for every open cover of $F$. 
        \end{proof}
        \item[5.] 
        \begin{proof}
        First of all the Cantor set is closed. Indeed the Cantor set is what's left over after removing all open intervals corresponding to the middle third of all segments that are used to construct the Cantor set. It follows that the complement of the Cantor set is the union of all those open intervals, which is open since arbitrary union of open sets is open. Therefore, by definition, the Cantor set is closed. Since the Cantor point is precisely the points $x \in [0,1]$ such that $x$ has a ternary expansion $x = 0.c_1c_2c_3c_4c_5 \cdots (\text{base 3}) = \sum_{k \in \NN, k = 1}^{\infty} \frac{c_k}{3^k}$ with $c_k \in \set{0,2}$. That means we can define a sequence $ a_n = \sum_{k \in \NN, k = 1}^{n} \frac{c_k}{3}$, taking the limit to infinity we recover the ternary expansion, $\lim_{n \to \infty} a_n = \sum_{k \in \NN, k = 1}^{\infty} \frac{c_k}{3^k}$, therefore $a_n \to x$ as $n \to \infty$ , so there exists a sequence $a_n$ that converges to $x$. It follows that $x$ is a cluster point. Since $x$ is arbitrary in $C$ , it follows that all points $x \in C$ are cluster points. In conclusion since $C$ is closed and all points are cluster points, $C$ is a perfect set. 
        \end{proof}
        \item[6.]
        \begin{proof}
        \textbf{lemma 1: any non empty open set in $\RR$ is a countable union of intervals with rational endpoints }\\
        First of all, since $\QQ$ is countable, the set $\mathcal{R}$ of all rational neighborhoods with rational radius centered around any $q \in \QQ$ countable. In other words, it's the infinite countable set of all open intervals with rational end points since $B_{r} (s ) = (r-s,r+s) = (q,p)$ with $r,s,q,p \in \QQ$. Second, recall that all open sets are union of open balls, since any open set $U \subseteq \RR$ has the property that $\forall x \in U$, $\exists \varepsilon > 0$ s.t $B_\varepsilon(x) \subset U$ and this is if and only if $U$ is the union of all such open balls. We can also say that any open set $U \subseteq \RR$ is the union of open intervals with real endpoints $(a,b)$, $a,b \in \RR$. Third, recall that, by the denseness of $\QQ$ in $\RR$, we can approximate any real number with a rational sequence. In other words, for any $r \in \RR$, there exists a $(q_k) \subseteq \QQ$ such that $q_k \to r$ as $k \to \infty$. Since by the denseness of $\QQ$ in $\RR$, $\exists q_k \in B_{\tfrac{1}{k}} (r)$, in fact there exists infinitely many rationals in $B_{\tfrac{1}{k}} (r)$, iteratively choosing $q_{k+1}$ such that $q_{k+1} \in B_{\tfrac{1}{k+1}} (r)$, we have that $\forall \varepsilon > 0 \quad  \exists K \in \NN$ s.t $\forall k \geq K$,  $\text{d} (q_k, r) \leq \text{d}(\frac{1}{k} , r) < \varepsilon$, so $q_k \to r $ as $k \to \infty$. And since $r$ is arbitrary, we can conclude that for any $r \in \RR$ there exists a sequence of rationals that converge to $r$. \textbf{Note that we can approach $r$ with a sequence of rationals either from the left or from the right such that a rational sequence $(q_k) \subseteq \QQ$ converges to $r$ with $q_k < r $ or $q_k > r$ for all $k$.}, It follows that for any open interval with real endpoints $(a,b)$ there exists a $(p_k) \subseteq \QQ$ such that $p_k \to a$ and a $(q_k) \subseteq \QQ$ such that $q_k \to b$, so $(p_k, q_k) \to (a,b)$ with $p_k > a$ and $q_k < b$ for all $k$. And in fact $\displaystyle \bigcup_{k = 1}^{\infty} (p_k, q_k)  = (a,b)$. We prove this equality. If $x \in \bigcup_{k = 1}^{\infty} (p_k, q_k)$ then $x \in (a,b)$ since starting from integer $k$, $a< p_k < x < q_k < b$. Now if $x \in (a,b)$, we need to show that $x$ is contained, for some integer $k$ in the interval $(p_k, q_k)$. Intuitively, since $p_k$ gets arbitrarily close to a, and $x$ is a fixed distance away from a, eventually $p_k < x$. Precisely, since $\forall k$ $a<p_k$ and $p_k \to a $, then $\forall \varepsilon < 0 $, $\exists K \in \NN$ such that $\forall k \geq K$, $p_k - a < \varepsilon$, choosing $\varepsilon = x-a$, we have $p_k - a < p_k - x \implies p_k < x$. Making the necessary changes, the proof is the same for showing that eventually $x < q_k$. In conclusion,  $\displaystyle \bigcup_{k = 1}^{\infty} (p_k, q_k)  = (a,b)$. And since every open interval, from analysis 1, is a disjoint,countable, union of open intervals, every $U_\alpha$ is a countable union of a countable union of rational open intervals. Thus, since every non empty open interval in $\RR$ is a union of members of $\mathcal{R}$, any non empty open set in $\RR$ is also such a union. \\
        \textbf{lemma 2: $\bigcup_{\alpha \in I}U_\alpha$ is exactly a countable union of rational intervals }\\
        Now, let $R = \set{R \in \mathcal{R} \ \colon \ \exists U \in \set{U_\alpha}_{\alpha \in I} \text{ with } R \subseteq U}$ be the countable set of all rational neighborhoods that are contained in some $U_\alpha$. Observe that this the union of all elements in this set is equal to $U_{\alpha \in I}U_\alpha$, listing out all the elements in $R = \set{R_1,R_2,...}$\\
        \textbf{claim:} $\bigcup_{R_i \in R} R_i = \bigcup \set{U_\alpha}_{\alpha \in I} $\\
        \begin{align*}
            x \in \bigcup_{R_i \in R} R_i &\implies x \in R_i \subseteq U_\alpha && \text{for some $i$ and $\alpha$} \\
            &\implies x \in \bigcup \set{U_\alpha}_{\alpha \in I} \\ 
        \end{align*}
        \begin{align*}
            x \in \bigcup \set{U_\alpha}_{\alpha \in I} &\implies x \in U_\alpha && \text{for some $\alpha$}\\
            &\implies x \in \bigcup_{R_i \in R} R_i && \text{ by lemma 1}\\
        \end{align*}
        Listing out the elements in $R = \set{R_1,R_2,R_3,\cdots }$ with each $R_i$ being a rational neighborhood, consider the set $U_0 = \set{ U_i \in \set{U_\alpha}_{\alpha \in I} \ \colon \ R_i \subseteq U_i \land i \in \NN \land R_i \in R}$. 
        \begin{align*}
         \overbrace{\bigcup U_0 \supseteq \bigcup R}^{\text{by construction}} = \bigcup \set{U_\alpha}_{\alpha \in I} \supset S \\
        \end{align*}
        \end{proof}
        so $U_0$ is a countable subcover
        \item[7.]
        \begin{itemize}
        	\item[(a)] Suppose towards a contradiction that every point $\alpha$ in the uncountable set $S$ has a neighborhood $r_\alpha$ such that $B_{r_0} (\alpha) \cap S$ is countable (and consequently every neighborhood smaller only has countable elements). We can now construct an open cover of $S$, which is the union of all countable neighborhoods of every $\alpha \in S$ 
        	\begin{align*}
        	S \subset \bigcup_{\alpha \in S} B_{r_\alpha} (\alpha) 
        	\end{align*}
        	But by question 6, every cover has a countable subcover, thus there exists $\set{\alpha_1,\alpha_2,\ldots} \subset S$ such that 
        	\begin{align*}
        		S \subset \bigcup_{\alpha \in \set{\alpha_1,\alpha_2,\ldots}} (B_{r_\alpha} (\alpha) \cap S) \subseteq \bigcup_{\alpha \in \set{\alpha_1,\alpha_2,\ldots}} B_{r_\alpha} (\alpha) 
        	\end{align*}
        	And this is a countable union of countable sets, which is countable, and since $S$ is no bigger than that union, $S$ is countable, which leads to a contradiction since $S$ is, by assumption, uncountable. 
        	\item[(b)]
        	Let $S_p$ be the set of all condensation points, let $(c_n) \subseteq S_p$ be a sequence of condensation points of $S$ and $(c_n) \to c$, we need to show that $c$ is a condensation point. We need to show that $\forall r > 0$ ,$ B_r(c) \cap S$ is uncountable. Since $c_n$ gets arbitrarily close to $c$ as n gets large, and $c_n$ is a condensation point for all $n$, then, we essentially need to show that there is uncountable set included in every neighborhood of $c$ by showing that there exists a neighborhood of a $c_n$ contained within every neighborhood of $c$. Indeed  $\forall \varepsilon > 0, \exists c_n \in B_\varepsilon(c)$ for some $n$ then take, $r = \frac{\dist{c_n,c}}{2}$, such that $B_r(c_n) \subset B_\varepsilon(c)$, and since $B_r(c_n) \cap S$, an uncountable set, has no less cardinality than $B_\varepsilon (c)\cap S$, So since $B_r(c_n) \cap S$ is uncountable, $B_\varepsilon(c)\cap S$ is uncountable. And since $\varepsilon$ is arbitrary, $c$ is a condensation point and thus in $S_p$, and we can conclude that the latter is closed. 
        	\item[(c)]
        	First of all, Since $\forall \varepsilon > 0 $ if $x$ is a condensation point, then $(B_\varepsilon(x) \cap S)\setminus \set{x}$ is uncountable, since removing one element from an uncountable set does not affect it's countability. Therefore, by (a), there exists a $x^\prime \in (B_\varepsilon(x) \cap S)\setminus \set{x}$ where $x \neq x^\prime$ by construction where $x^\prime$ is a condensation point. Furthermore, keeping in mind that removing a countable set of points from an uncountable set does not alter it's countability as well, we define:
        	\begin{align*}
        	 c_1 &\in (B_1(x)\cap S) \setminus \set{x} \\
        	 c_2 &\in (B_\frac{1}{2}(x) \cap S) \setminus \set{x,c_1} \\
        	 c_3 &\in (B_\frac{1}{3}(x) \cap S) \setminus \set{x,c_1,c_2} \\
        	 &\vdots\\
        	 c_k &\in (B_\frac{1}{k} (x) \cap S) \setminus \set{x,c_1,c_2,\ldots,c_{k+1}}
        	\end{align*}
        	to be points in the sequence $c_k$ of condensation points of $S$, then since $\forall \varepsilon > 0, \exists K \in \NN$ such that $\forall k \geq K$, $\dist{c_k,x} < \frac{1}{k} < \varepsilon$, then $c_k \to x$, and since $x \in S$ is arbitrary, every $x \in S$ is a cluster point, so $S$ is perfect and has no isolated points. 
        	\item[(d)]
        	Let $P$ be the set of all points in $S$ that are not condensation points. That means if $x \in P$, then there exists a non empty neighborhood $V_x$ of $x$ such that $V \cap S$ is countable. Additionally, observe that 
        	\begin{align*}
        	 \bigcup_{x \in P} V_x
        	\end{align*}
        	is an open over for $P$. But by question 6, every open cover has a countable subcover, so there exists $\set{x_1,x_2,\ldots}$ such that 
        	\begin{align*}
        	P \subset \bigcup_{x \in \set{x_1,x_2,\ldots}} V_x
        	\end{align*}
        	is an open cover of $P$, and since it's a countable union of countable sets, it's countable, and since the cardinality of $P$ is no greater than $\displaystyle \bigcup_{x \in \set{x_1,x_2,\ldots}} V_x$, $P$ is countable.
        	\item[8.]
        	We claim that a non empty closed set is the union of the set of it's condensation points that are in $S$, call it $S_c$ with the set of points in $S$ that are not condensation points. Observe that 
         \begin{align*}
             S = S \cap (S_c \cup S_c^c)  \\
             &= \underbrace{(S \cap S_c)}_{\text{closed under intersection}} \cup (S \cup S_c^c) \\
             &= S_c \cup (S \cap S_c^c)
         \end{align*}
         $S \cap S_c^c$ is countable since $S_c^c$ is countable by the previous question. Since $S$ is closed and $S_c$ is closed, then their intersection is closed. Additionally, $S_c \cap S = S_c$, since any point in $S_c$ is in $S$ if $S$ is closed. This follows from that fact that $S$ has to contain all it's limit points, and since $S_c$ is the set of all condensation points, every point in $S_c$ is trivially a limit point, since being a condensation means that in every neighborhood there exists an uncountably infinite amount of points in $S$, while a limit point just needs a infinite amount of points in every neighborhood. And $S_c$, by the previous question is perfect, so in conclusion, S is the union of a perfect set and an uncountable set. 
        	\item[9.]
        	Let $\varepsilon > 0 $, then by continuity of $f$, $\forall x \in [a,b]$, $\exists \delta_x > 0 $ such that $f(B(x,\delta_x)) \subseteq B(f(x), \frac{\varepsilon}{2})$. Notice that all of the $\delta_x$ form an open cover for $[a,b]$, but so does all the $\set{B(x_,\frac{\delta_x}{2})}$. So since $[a,b]$ is compact by Heine-Borel, then we can find a finite subcover $\set{B(x_i, \frac{\delta_i}{2})}_{i = 1,2,3,...,n}$. We need to show that there exists a $\delta > 0 $ such that for any two points $x,y \in [a,b]$, if they are within $\delta$ of each other, that means $|x-y| < \delta$, then $\abs{x-y} < \delta_i$ for some $i$. Then we can just choose $\delta$ to the minimum  $\delta = \min \set{\delta_1/2,\delta_2/2,...,\delta_n/2}$, we can do this since it's a finite set with all $\delta_i > 0$. We can let $x \in [a,b]$ and it will for sure be in $B(x_i, \frac{\delta_i}{2})$ for some $i$ since it covers the entirely of $[a,b]$, then if $\abs{x-y} < \delta$ that means, by the triangle inequality: $\abs{x_i-y} \leq \abs{x_i-x} + \abs{x-y} < \frac{\delta_i}{2} + \min \set{\delta_1/2,\delta_2/2,...,\delta_n/2} \leq \frac{\delta_i}{2}  + \frac{\delta_i}{2} = \delta_i $, so $y \in B(x_i, \delta_i)$, and by definition $f(B(x_i,\delta_i)) \subseteq B(x_i,\frac{\varepsilon}{2})$, and so $\abs{f(x)-f(y)} \leq \abs{f(x) - f(x_i)} + \abs{f(y) - f(x_i)} < \frac{\varepsilon}{2} + \frac{\varepsilon}{2} = \varepsilon$, $f$ is also uniformly continuous on $[a.b]$.
        \end{itemize}
    \end{itemize}
    \end{document}
